\chapter{結論}
第1〜3回の実験から、同一validation上のMetric Outdoor Smallは未学習RMSE14.5340 mから1 epoch適応Cの1.7490 m、独立5 epoch run Dのepoch 5で1.0094 mへ改善した。D epoch 5のMAEは0.7373 m、AbsRelは0.03694、Pearsonは0.9919で、絶対深度誤差と全14 px局所差MAEは低下傾向にあった。epoch 2では絶対誤差が一時増加し、複数指標が同じ方向に単調改善するとは限らなかった。
局所高低差には課題が残る。CとD epoch 5のGT差20 cm以上の振幅比は報告値でいずれも約0.33であり、大きな差分の過小推定が残った。ただしrun間の詳細標本抽出条件を完全には照合できず、この近似値の小差を定量的な改善効果とは解釈しない。第2回ではRelative由来Aの局所誤差が画像端128 pxへ偏り、内部でモデル順位が変わった。ゼロ差分基準とGT階級分析は、総局所MAEが平坦予測に有利になり得ることを示した。
確認済みなのは4作物のDay047 validationにおけるZ-depth結果である。稲垂直草丈、未知撮影日汎化、independent test、20 epoch完了は未検証である。次段階では保存checkpointから基準runを継続評価し、別runの局所損失・高解像tile比較、地上実測高さを含む稲評価へ進む必要がある。
